\chapter{生物信息词典}

\section*{这本词典怎么用}

读生物学新闻或课本时，你会反复遇到这十五个词：DNA、RNA、基因、等位基因、染色体、基因组、转录组、蛋白质组、基因型、表型、遗传率、突变、变异、表达、调控。它们看起来都像“遗传学黑话”，其实描述的是同一条信息流的不同环节：信息\textbf{存在哪里}、\textbf{以什么为单位}、\textbf{怎样被拿出来用}、\textbf{最终变成什么性状}。

所以本词典不按拼音排序，而是按这条信息流分成四组：

\begin{itemize}[itemsep=2pt]
\item \textbf{第一组　信息的载体}：DNA、RNA、染色体、基因组——信息写在什么材料上，怎样打包。
\item \textbf{第二组　信息的单位与版本}：基因、等位基因、突变、变异——信息怎样分段，版本怎样产生差异。
\item \textbf{第三组　信息的执行}：表达、调控、转录组、蛋白质组——信息怎样被读取、被调节、被清点。
\item \textbf{第四组　信息与性状}：基因型、表型、遗传率——信息和你能看见的性状之间，到底隔着多少东西。
\end{itemize}

每个词条分两部分：先用一段“人话”说清楚它是什么，再专门指出它最容易和谁混淆、最容易被怎样误解。词条之间互相引用，括号里注明全书正文中展开讲述的章节，方便你顺藤摸瓜。如果你只想花三十秒回忆某个词，直接翻最后的速查表；如果你发现自己把两个词混着用，请回到对应词条把“容易和谁混淆”那段重读一遍——这本词典一半的价值就藏在那些辨析里。查词典只是第一步，真正的理解在正文的因果链里。

\section{第一组　信息的载体}

\subsection{DNA（脱氧核糖核酸）}

\textbf{一句话定义：} DNA 是一种超长的链状分子，生命把遗传信息写在这条链的“字母顺序”里。

它的字母只有四种：A、T、G、C（四种碱基）。一条 DNA 链就是这四个字母排出的一长串序列，比如 …A-T-G-C-C-A-T…。细胞里的 DNA 通常是两条链并排拧成的双螺旋，两条链按“A 对 T、G 对 C”的规则互补配对——知道一条链，就能推出另一条。这正是它能被忠实抄写（复制）的结构基础，第 66 章专门讲了“为什么这个结构天生适合保存信息”，第 65 章讲了科学家怎样确认遗传物质是 DNA 而不是蛋白质。你的每个细胞核里，DNA 拉直了大约有 2 米长，写着约 30 亿个字母。

\textbf{容易和谁混淆：} DNA \textbf{不等于}基因。DNA 是整张纸，基因只是纸上某些有意义的段落（见“基因”条）——你的 DNA 里只有一小部分是基因。更要警惕一句流行语：“DNA 是决定你的一整套程序代码。”这是夸大。DNA 更像一套会被环境反复批注的食谱：做什么菜受它影响，但火候、食材、厨师的手艺都会改变端上桌的结果。从一段 DNA 到一个真实的你，中间隔着整个生命过程（第 73 章）。

\subsection{RNA（核糖核酸）}

\textbf{一句话定义：} RNA 是 DNA 的“近亲”分子，通常是单链，负责把信息从 DNA 搬运、转换成细胞能用的形式。

它和 DNA 的差别主要有三处：糖环上多一个氧原子；通常是单链而不是双链；字母表里用 U 代替 T。但别只把它当“信使”——RNA 在细胞里至少干三种活：mRNA 抄写基因的内容送往蛋白质工厂（第 68 章），tRNA 在工厂里当“翻译员”核对氨基酸（第 69 章），rRNA 干脆就是蛋白质工厂（核糖体）的骨架。有些病毒干脆用 RNA 当自己的遗传物质，比如流感病毒。

\textbf{容易和谁混淆：} RNA 不是“残缺版 DNA”或“DNA 的临时工”。它是一类功能独立的分子，甚至很可能是生命史上先出现的那位主角——既能存信息又能催化反应，今天的 DNA—RNA—蛋白质分工也许是后来才演化出来的（生命起源的细节仍有不少争论，第 87 章会坦白哪些还不知道）。另外，mRNA 疫苗里的“mRNA”就是它，作用只是给细胞递一份临时说明书，几天内就被细胞正常降解，不会改写你的 DNA。

\subsection{染色体}

\textbf{一句话定义：} 染色体是 DNA 缠绕蛋白质后打包成的“线轴”，是 DNA 在细胞里的运输形态。

2 米长的 DNA 要塞进直径只有头发丝几十分之一的细胞核，靠的是层层缠绕：DNA 绕在一团团叫组蛋白的蛋白质上，像线绕在线轴上，线轴再折叠、再压缩。平时它以松散的“染色质”状态铺开，方便细胞读取；到细胞分裂时，才压缩成显微镜下可见的粗短小棍。人有 46 条染色体，配成 23 对——其中 23 条来自父亲，23 条来自母亲（第 64 章）。

\textbf{容易和谁混淆：} 染色体和基因不是一个层面的东西：一条染色体是一条超长的 DNA 加蛋白质，上面排着成百上千个基因。教科书里那个“X 形”的染色体图像也误导了很多人——那只是细胞分裂中期、每条 DNA 已经复制成两份（两条姐妹染色单体）时的临时造型，不是它平时的样子。还有一点：染色体的“对”是成对携带同一套基因的不同版本，不是“两条完全相同的拷贝”（见“等位基因”条）。最后提醒一句：染色体数目异常（比如 21 号染色体多了一条，即唐氏综合征）是“整套打包”层面的问题，和某个基因突变是不同层级的事故，新闻里常把两者混着讲。

\subsection{基因组}

\textbf{一句话定义：} 基因组指一个细胞（或一个物种）所携带的全部遗传信息的总和。

如果说 DNA 是纸、基因是段落，基因组就是整套藏书——包括所有基因，也包括基因之间大量不编码蛋白质的序列。人类基因组约 30 亿个碱基对，其中真正编码蛋白质的部分只占大约百分之一二，其余的不是“垃圾”：很多段落参与调控、参与染色体的结构维护，也有不少是演化留下的“旧文件”（第 79 章讲了基因组怎样复制、借用和改造自己）。2003 年人类基因组计划完成第一版完整测序；今天测一个人的全基因组只要几百美元和几天时间（第 81 章）。

\textbf{容易和谁混淆：} 基因组 \textbf{不等于}“所有基因的集合”——它包含基因，也包含基因以外的一切序列，这正是“组”字的意思：全部家底一起算。也要区分“一个物种的基因组”和“某个人的基因组”：人类有一个参考基因组，但你和同学的基因组在约千分之一的字母上互不相同——那一点点差异，加上环境，就是你和你的同桌不一样的原因之一。

\section{第二组　信息的单位与版本}

\subsection{基因}

\textbf{一句话定义：} 基因是 DNA 上能产生一个有用产物（通常是某种蛋白质，有时是某种 RNA）的一段序列。

它是遗传信息的基本“条目”：一段有起点、有终点、能被细胞读取使用的文字。人大约有两万个蛋白质编码基因。这个数字曾让科学家大跌眼镜——人类基因组计划启动前，很多人打赌至少有十万个，毕竟水蚤都有约三万个。基因的数量显然不决定生物的复杂程度，真正的差别在信息的使用方式上（见“调控”条）。基因的概念其实很年轻：孟德尔 1866 年只算出“有某种遗传因子”（第 63 章），“基因”这个词 1909 年才被造出来，而“基因到底是什么”这个问题，科学家到今天还在细化答案（第 71 章专门讨论）。

\textbf{容易和谁混淆：} 第一，基因不是串在染色体上的一颗颗独立的珠子——有些基因的地盘互相重叠，一段 DNA 还能通过不同的读法产生好几种产物。第二，不存在“完美基因”或“坏基因”的绝对名单：同一个基因版本，在一种环境里帮倒忙，在另一种环境里可能救命——镰状细胞贫血的那个突变版本，在疟疾流行地区就保护过携带者（第 72、76 章）。第三，“有一个基因”和“这个基因正在工作”是两回事（见“表达”条）。

\subsection{等位基因}

\textbf{一句话定义：} 等位基因是同一个基因的不同版本，就像同一款软件的不同发行版。

你的 23 对染色体里，每一对的两条在相同位置携带同一批基因，但具体版本可能不同：决定 ABO 血型的基因，有人是 A 版，有人是 B 版，有人是 O 版。这些版本就是等位基因。你从父亲和母亲那里各得到一个版本；有些版本之间是“显性—隐性”关系（杂合时只表现出显性版本的效果），有些是共同表现（A 版和 B 版凑成 AB 型血）。

\textbf{容易和谁混淆：} 这是对误解重灾区。显性等位基因\textbf{不是}更强、更常见、更有利的版本——“显性”只描述杂合时哪个版本的效果在表型上显现，跟好坏、多少毫无关系。隐性版本没有消失，它还在 DNA 里好好地存着，照样能传给下一代，遇到另一个隐性版本时就会重新显现。O 型血是隐性，却是很多地区最常见的血型；引起苯丙酮尿症的等位基因是隐性，可它一点都不“弱”。另外，等位基因说的是同一个基因内部的版本差异，不要把“不同的基因”也叫成等位基因。

\subsection{突变}

\textbf{一句话定义：} 突变是 DNA 序列发生的任何改变——字母写错、漏掉、多出或整段重排。

它怎么来？最常见的是复制时的抄写错误：细胞每分裂一次要抄 30 亿个字母，校对系统再强也会漏网（第 67 章）；也可能来自紫外线、某些化学物质对 DNA 的损伤没修好。突变有大有小：换错一个字母可能毫无影响，也可能改变一个氨基酸从而改变整个蛋白质——镰状细胞贫血就只是血红蛋白基因里一个字母的差别。

\textbf{容易和谁混淆：} 两个方向都想歪。一种把突变全当成灾难：实际上绝大多数突变是中性的，既没好处也没坏处；少数有害；极少数在特定环境里有用——而进化恰恰靠这极少数积累出新功能（第 72、76 章）。另一种把突变化身成“有求必应”：细胞不会因为“需要”某个能力就定向产生对应突变。突变是随机发生的拼写事故，自然选择只是事后筛子——先有了五花八门的版本，环境再决定哪些版本留下。顺序不能反。

\subsection{变异}

\textbf{一句话定义：} 变异指同种生物个体之间存在的差异——可以是 DNA 层面的，也可以是性状层面的。

班里同学身高不同、血型不同、对咖啡因代谢快慢不同，这些都是变异。没有变异，自然选择就没有原材料，进化就停摆；可以说，变异是进化的燃料（第 76、77 章）。变异一部分来自可遗传的 DNA 差异（突变的积累、父母版本的重新组合），一部分来自环境造成的差别（营养、锻炼、经历），很多时候两者纠缠在一起。

\textbf{容易和谁混淆：} 变异和突变最容易被当成同义词，其实层级不同：突变是一次“产生新序列的事件”，变异是“群体中差异存在的状态”。突变是制造变异的来源之一，但变异还有别的来源——比如有性生殖把父母已有的版本重新洗牌，不产生任何新突变也能制造海量新组合。还要注意：变异不“朝某个方向”。群体里的差异是散开的，没有目标；朝着环境走的那一步是选择，不是变异自己长出了方向。

\section{第三组　信息的执行}

\subsection{表达}

\textbf{一句话定义：} 基因表达，就是细胞把某个基因的信息“拿出来用”的过程——先抄成 RNA（转录，第 68 章），再按需翻译成蛋白质（第 69 章）。

你的每个细胞都装着同一套基因组，但肝细胞在表达解毒相关的基因，神经细胞在表达传递信号相关的基因，皮肤细胞在表达角蛋白相关的基因。“表达”还有强弱的程度之分：同一个基因，可以大量表达，也可以只零星表达，就像音响的音量旋钮，而不是只有开关两档。你饿的时候和饱的时候，身体里同一批基因的表达量都在变化（第 70 章）。

\textbf{容易和谁混淆：} “有这个基因”不等于“这个基因正在表达”。基因检测报告说“你携带某基因”，说的只是序列在；它有没有被使用、用了多少、在哪些细胞里用，是另一回事。日常语言里“表达情绪”“表达观点”的“表达”是“表现出来”的笼统意思，生物学里的“基因表达”有严格所指：从 DNA 到 RNA 再到蛋白质这条具体的流水线。也别以为基因表达是一锤子买卖——它是一个随时可以调高、调低、暂停的持续过程。

\subsection{调控}

\textbf{一句话定义：} 调控是决定“哪个基因、在什么时间、什么地点、表达多少”的全套管理机制。

基因组里的两万个基因如果全部同时开工，细胞立刻乱成一锅粥。所以每个基因旁边都配有“控制面板”：启动子等序列是面板本身，转录因子等蛋白质是按面板的手。细胞外的一餐饭、一次感染、一天的光照变化，都能通过信号链条最终按到这些面板上。还有一层更“持久”的调控：不改变 DNA 序列，只给 DNA 或它的包装蛋白做化学标记，让某些基因长期静音——这就是表观遗传的领域（第 74 章）。受精卵能从一个细胞发育出两百多种细胞，靠的几乎全是调控的分工，而不是基因组的差异（第 75 章）。

\textbf{容易和谁混淆：} 调控不改变基因的序列，它改变的是序列的“使用方式”——这一点和突变正好相反：突变改字，调控改读法。很多疾病的根源不在基因“坏了”，而在调控“乱了”：基因本身完好，却在错误的时间、错误的组织里被打开或关掉，不少癌症就是这样起步的。所以看到“某某疾病是基因导致的”，正确的追问是：是序列变了，还是使用方式变了？两者都很常见，含义完全不同。

\subsection{转录组}

\textbf{一句话定义：} 转录组是某个细胞（或组织）在某一时刻全部被转录出来的 RNA 的总集合。

如果说基因组是整套藏书，转录组就是“此刻被借出来复印的那些页面”的清单。这份清单是高度动态的：你的肝细胞和神经细胞共享同一个基因组，转录组却大不相同——这正是它们长得不一样、干的活不一样的原因。同一个细胞在发烧前后、饭前饭后，转录组也会改写。科学家研究“癌细胞和正常细胞差在哪”“这株抗旱水稻打开了哪些基因”，测的往往就是转录组（第 81 章讲了这类测量技术背后的测序原理）。

\textbf{容易和谁混淆：} 转录组最大的对照对象是基因组。基因组相对稳定——你从受精卵到老，每个细胞里的基因组基本一样；转录组却随时随地、因细胞类型而异。体检式的“基因检测”测的是基因组（相对稳定的那套序列），而判断“某个组织此刻在干什么”需要的是转录组。别把两者混为一谈：一个告诉你“图书馆里有什么书”，另一个告诉你“此时此刻谁在读哪几本”。

\subsection{蛋白质组}

\textbf{一句话定义：} 蛋白质组是某个细胞（或组织）在某一时刻存在的全部蛋白质的总集合。

蛋白质才是真正干活的分子：催化反应的酶、运输氧气的血红蛋白、收缩肌肉的肌动蛋白、传递信号的受体——生命的大多数“动作”都是蛋白质做的（第 44 到 49 章讲了它们为什么能干这么多活）。蛋白质组比转录组还要多变：蛋白质合成后还能被贴上磷酸、糖链等“便利贴”修饰，被切割、被降解，数量和活性时刻在调整。

\textbf{容易和谁混淆：} “测了转录组就等于知道蛋白质组”是常见错觉。RNA 多不代表蛋白质一定多：从 RNA 到蛋白质还有翻译效率这一关，蛋白质还有各自的寿命长短。同一份 mRNA，在不同条件下可能产出差别好几倍的蛋白质。还要记住蛋白质组不等于基因组：基因组回答“能造什么”，转录组回答“正在抄写哪些图纸”，蛋白质组回答“此刻车间里有哪些机器在转”——三个“组”是同一信息流的三个清点站，一个比一个接近“正在发生的生命”。

\section{第四组　信息与性状}

\subsection{基因型}

\textbf{一句话定义：} 基因型指一个生物的遗传组成——可以指整套基因组，也常指某个具体基因上携带的等位基因组合。

说“这个人的 ABO 基因型是 AO”，意思是：在血型这个基因的两个位子上，他从父母那里拿到了一个 A 版和一个 O 版。基因型是写在 DNA 里、相对稳定、可以传给下一代的那部分信息。它是“图纸”层面的事实，不是看得见摸得着的特征。

\textbf{容易和谁混淆：} 基因型和表型是一对必须成组记的词（见“表型”条）：基因型是图纸，表型是盖出来的房子加装修。最容易犯的错是从表型直接反推基因型——看到 A 型血，就说“基因型是 AA”，其实 AO 也表现为 A 型血（因为 A 对 O 显性）。同一个表型背后可能有不同基因型，反过来，同一基因型在不同环境里也会长出不同的表型。基因检测报告上那些字母组合是基因型层面的信息，它能提示概率和倾向，却不能直接朗读出“你会成为什么样的人”。

\subsection{表型}

\textbf{一句话定义：} 表型是一个生物所有能观察、能测量的性状的总和——从身高、血型到代谢速度、对某种药的反应。

关键公式只有一句话：\textbf{表型 $=$ 基因型 $\times$ 环境}（更诚实地说，还要加上一点发育中的随机噪声）。同卵双胞胎基因型几乎相同，分开抚养仍会长出细微差别的身高、体重和疾病史——那就是环境和随机性的签名。反过来，同样“每天跑步五公里”的两个人，肌肉增长幅度可能差很多——那是基因型的签名。再举一个田间的例子：同一批种子分种在肥沃和贫瘠的两块地里，株高立刻分出两档——基因型没变，表型变了。从一段 DNA 到一个性状，中间隔着 RNA、蛋白质、细胞、组织、激素、营养和运气（第 73 章整章都在拆这条长链）。

\textbf{容易和谁混淆：} 表型不只是“长相”。血型、乳糖耐受与否、某种酶的活性、患病风险，凡是能测量的性状都算表型——很多重要的表型肉眼根本看不见。另一个误区是以为表型里“基因的部分”和“环境的部分”可以像切蛋糕一样分开到个人头上：对单个个体问“你的身高几分之几来自基因”是没有意义的问题——只有在群体层面谈差异来源才有意义（见“遗传率”条）。

\subsection{遗传率}

\textbf{一句话定义：} 遗传率是一个统计量：在某个特定群体、特定环境里，个体间某性状的\textbf{差异}有多大比例能用遗传差异来解释。

注意它描述的是“差异的差异”，取值在 0 到 1 之间。比如身高在营养普遍充足的人群里遗传率可高达 0.8 左右——意思是：这群人高矮不齐这件事，大约八成能与遗传差异相关。它靠双胞胎研究、收养研究等方法估算，第 73 章讲了这类研究的设计和它们排除不了什么。

\textbf{容易和谁混淆：} 这是本词典里被误解最狠的词，请慢读三遍：遗传率\textbf{不是}“一个人的性状有多少百分比由基因决定”。“身高遗传率 0.8”绝不等于“你的身高 80\% 来自基因、20\% 来自营养”——对一个具体的人，这个问题根本不成立，就像问“这场考试的分数，几分之几来自试卷”一样没有意义。遗传率还依附于群体和环境：如果让所有人吃到完全一样的饭，身高差异中环境的份额缩小，遗传率的数字反而会升高——同一个性状，换个人群、换个年代，遗传率就变了。所以“遗传率高”也不等于“没法改变”：苯丙酮尿症由基因引起、遗传率极高，但只要出生后严格控制饮食中的苯丙氨酸，患儿可以正常发育——环境的干预照样有效。下次看到“某性状百分之几十由基因决定”的标题，你就知道该怎么皱眉了（附录 E 会继续教你拆解这类新闻）。

\section*{一张速查表}

\begin{table}[htbp]
\centering
\begin{tabular}{llp{6.2cm}}
\toprule
\textbf{词条} & \textbf{属于} & \textbf{一句话提醒} \\
\midrule
DNA & 载体 & 写字母的材料，不等于基因 \\
RNA & 载体 & 多面手，不只是信使 \\
染色体 & 载体 & DNA 的打包运输形态 \\
基因组 & 载体 & 全部遗传信息，含非编码序列 \\
基因 & 单位 & 能产生有用产物的一段 DNA \\
等位基因 & 版本 & 同一基因的不同版本，显性不等于更好 \\
突变 & 版本 & 序列的改变，多数中性，随机发生 \\
变异 & 版本 & 群体中的差异，进化的原材料 \\
表达 & 执行 & 把基因拿出来用，有强有弱 \\
调控 & 执行 & 决定何时何地用多少，不改序列 \\
转录组 & 执行 & 此刻抄出来的全部 RNA \\
蛋白质组 & 执行 & 此刻干活的全部蛋白质 \\
基因型 & 性状 & 图纸层面：携带哪些版本 \\
表型 & 性状 & 可观测的性状 $=$ 基因型 $\times$ 环境 \\
遗传率 & 性状 & 群体差异的统计量，不能用于个人 \\
\bottomrule
\end{tabular}
\end{table}

\section*{最后一句提醒}

这十五个词其实只讲了一件事：信息怎样从一段稳定的 DNA 序列，经过层层使用、调节和环境的批注，最终变成一个活生生的、会变化的个体。记住这条主线，词典里的每个词条都能挂回它自己的位置；忘了这条主线，它们就会退化成十五个孤立的考点——而这正是这本书想帮你烧掉的那种学法。正文第 65 到 75 章是这条主线的完整旅程，遇到任何含糊的词条，请回到那里去。
